A final observation concerns how saliency behaves for \emph{incorrect} predictions. When a VLM miscounts---answering ``four'' for an image of three cats---we often (though not always) see a phantom blob of attention in the background, structurally similar to the mass on the three real objects, as if the model attends to evidence for an object that is not there. The complementary error also occurs: when the model undercounts (``two'' for three cats), the saliency skips the missed object entirely. The effect was visible in somewhat more than half the images we examined, most clearly for large objects and small counts. We report it as a qualitative phenomenon; quantifying this ``saliency hallucination'' and relating it to answer errors is an interesting direction we did not pursue here.
